\documentclass[12pt,a4paper]{ctexart}

% ============================================================
%                       宏包设置
% ============================================================

\usepackage{amsmath, amssymb, amsfonts}
\usepackage{geometry}
\usepackage{fancyhdr}
\usepackage{graphicx}
\usepackage{xcolor}
\usepackage{setspace}
\usepackage[most]{tcolorbox}
\usepackage{enumitem}

% ============================================================
%                 清新蓝绿色配色
% ============================================================

% 主色：清爽青绿色
\definecolor{SZURed}{HTML}{2F9E8F}

% 题目区域：浅薄荷绿
\definecolor{SZULightRed}{HTML}{EAF7F4}

% 答案区域：浅天空蓝
\definecolor{SZULightRose}{HTML}{EEF7FC}

% 边框：柔和蓝绿色
\definecolor{SZUSoftBorder}{HTML}{9CCFC8}

% 强调文字：沉稳蓝绿色
\definecolor{SZUTextRed}{HTML}{256B72}

% ============================================================
%                       页面设置
% ============================================================

\geometry{
    left=2.5cm,
    right=2.5cm,
    top=2.5cm,
    bottom=2.5cm
}

\setlength{\parindent}{0pt}
\setlength{\parskip}{0.6em}
\linespread{1.15}

% ============================================================
%                       课程信息
% ============================================================

\newcommand{\coursename}{概率论与数理统计}
\newcommand{\department}{深圳大学生物医学工程}
\newcommand{\homeworknumber}{作业 6}

% ============================================================
%                  学生信息 —— 请学生修改
% ============================================================

\newcommand{\studentname}{请填写姓名}
\newcommand{\studentid}{请填写学号}

% ============================================================
%                       页眉页脚
% ============================================================

\pagestyle{fancy}
\fancyhf{}

\lhead{\textcolor{SZUTextRed}{\coursename}}
\rhead{\textcolor{SZUTextRed}{\homeworknumber}}
\cfoot{\textcolor{SZUTextRed}{\thepage}}

\renewcommand{\headrulewidth}{0.4pt}
\renewcommand{\headrule}{
    \hbox to\headwidth{
        \color{SZUSoftBorder}\leaders\hrule height \headrulewidth\hfill
    }
}

\setlength{\headheight}{15pt}

% ============================================================
%                    题目编号与环境设置
% ============================================================

\newcounter{problem}

% -------------------------
% 题目环境
% -------------------------

\newenvironment{problem}[1][]
{
    \stepcounter{problem}
    \begin{tcolorbox}[
        title={题目 \theproblem\if&#1&\relax\else\quad[#1]\fi},
        colback=SZULightRed,
        colframe=SZURed,
        coltitle=white,
        fonttitle=\bfseries,
        boxrule=0.8pt,
        arc=2mm,
        left=3mm,
        right=3mm,
        top=2.5mm,
        bottom=2.5mm,
        breakable
    ]
}
{
    \end{tcolorbox}
}

% -------------------------
% 答案环境
% -------------------------

\newenvironment{answer}
{
    \begin{tcolorbox}[
        title={学生答案},
        colback=SZULightRose,
        colframe=SZUSoftBorder,
        coltitle=SZUTextRed,
        fonttitle=\bfseries,
        boxrule=0.7pt,
        arc=2mm,
        left=3mm,
        right=3mm,
        top=2.5mm,
        bottom=2.5mm,
        breakable
    ]
}
{
    \end{tcolorbox}
}

% ============================================================
%                         正文
% ============================================================

\begin{document}

% ============================================================
%                         标题
% ============================================================

\begin{center}

    {\LARGE\bfseries\textcolor{SZURed}{\coursename}}

    \vspace{0.4em}

    {\large \department}

    \vspace{1em}

    {\Large\bfseries\textcolor{SZUTextRed}{\homeworknumber}}

\end{center}

\vspace{1em}

% ============================================================
%                       学生信息
% ============================================================

\begin{center}

\begin{tabular}{rl@{\hspace{3cm}}rl}

\textcolor{SZUTextRed}{\bfseries 姓名：}
&
\studentname
&
\textcolor{SZUTextRed}{\bfseries 学号：}
&
\studentid

\end{tabular}

\end{center}

\vspace{0.7em}

{\color{SZUSoftBorder}\hrule}

\vspace{1.5em}



% ============================================================
% ============================================================
%
%                           题目 1
%
% ============================================================
% ============================================================

\begin{problem}[课本习题2(2,3)]

设 $X_1,X_2,\dots,X_n$ 为总体的一个样本，$x_1,x_2,\dots,x_n$ 为一相应的样本值。求下列各总体的概率密度或分布律中的未知参数的矩估计量和矩估计值。
\begin{enumerate}[label=(\arabic*)]
    \item[(2)]
\[
f(x)=
\begin{cases}
\sqrt{\theta}\,x^{\sqrt{\theta}-1}, & 0\le x\le 1,\\
0, & \text{其他},
\end{cases}
\]
其中 $\theta>0$，$\theta$ 为未知参数。
    \item[(3)]
\[
P\{X=x\}=
\binom{m}{x}p^{x}(1-p)^{m-x},\quad x=0,1,2,\dots,m,
\]
其中 $0<p<1$，$p$ 为未知参数。
\end{enumerate}
\end{problem}


% ============================================================
%
%              ↓↓↓↓↓↓ 学生答案区域 ↓↓↓↓↓↓
%
%          请只在 answer 环境中填写你的答案
%
% ============================================================

\begin{answer}

% ==========================================================
%                 请从下一行开始填写答案
% ==========================================================






% ==========================================================
%                 请在上一行结束填写答案
% ==========================================================

\end{answer}

% ============================================================
%
%              ↑↑↑↑↑↑ 学生答案区域 ↑↑↑↑↑↑
%
% ============================================================


\vspace{1em}



% ============================================================
% ============================================================
%
%                           题目 2
%
% ============================================================
% ============================================================

\begin{problem}[课本习题3]

求上题中各未知参数的最大似然估计值和估计量。
\end{problem}


% ============================================================
%
%              ↓↓↓↓↓↓ 学生答案区域 ↓↓↓↓↓↓
%
% ============================================================

\begin{answer}

% ==========================================================
%                 请从下一行开始填写答案
% ==========================================================





% ==========================================================
%                 请在上一行结束填写答案
% ==========================================================

\end{answer}

% ============================================================
%
%              ↑↑↑↑↑↑ 学生答案区域 ↑↑↑↑↑↑
%
% ============================================================


\vspace{1em}



% ============================================================
% ============================================================
%
%                           题目 3
%
% ============================================================
% ============================================================

\begin{problem}[课本习题10]

设 $X_1,X_2,\dots,X_n$ 是来自总体 $X$ 的一个样本，设 $E(X)=\mu$，$D(X)=\sigma^2$。
\begin{enumerate}[label=(\arabic*)]
    \item 确定常数 $c$，使 $c\sum_{i=1}^{n-1}(X_{i+1}-X_i)^2$ 为 $\sigma^2$ 的无偏估计。
    \item 确定常数 $c$ 使 $(\bar{X})^2-cS^2$ 是 $\mu^2$ 的无偏估计（$\bar{X},S^2$ 是样本均值和样本方差）。
\end{enumerate}
\end{problem}


% ============================================================
%
%              ↓↓↓↓↓↓ 学生答案区域 ↓↓↓↓↓↓
%
% ============================================================

\begin{answer}

% ==========================================================
%                 请从下一行开始填写答案
% ==========================================================





% ==========================================================
%                 请在上一行结束填写答案
% ==========================================================

\end{answer}

% ============================================================
%
%              ↑↑↑↑↑↑ 学生答案区域 ↑↑↑↑↑↑
%
% ============================================================


\vspace{1em}



% ============================================================
% ============================================================
%
%                           题目 4
%
% ============================================================
% ============================================================

\begin{problem}[课本习题11]

设总体 $X$ 的概率密度为
\[
f(x;\theta)=
\left\{
\begin{array}{llr}
\dfrac{1}{\theta}x^{(1-\theta)/\theta}, & 0<x<1, & \\
& & 0<\theta<+\infty,\\
0, & \text{其他}. &
\end{array}
\right.
\]
$X_1,X_2,\dots,X_n$ 是来自总体 $X$ 的样本。
\begin{enumerate}[label=(\arabic*)]
    \item 验证 $\theta$ 的最大似然估计量是 $\hat{\theta} = \dfrac{-1}{n}\sum_{i=1}^{n}\ln X_i$。
    \item 证明 $\hat{\theta}$ 是 $\theta$ 的无偏估计量。
\end{enumerate}
\end{problem}


% ============================================================
%
%              ↓↓↓↓↓↓ 学生答案区域 ↓↓↓↓↓↓
%
% ============================================================

\begin{answer}

% ==========================================================
%                 请从下一行开始填写答案
% ==========================================================





% ==========================================================
%                 请在上一行结束填写答案
% ==========================================================

\end{answer}

% ============================================================
%
%              ↑↑↑↑↑↑ 学生答案区域 ↑↑↑↑↑↑
%
% ============================================================





\end{document}
```
